%%%%%%%%%%%%%%%%%%%%%%%%%%%%%%%%%%%%%%%%%
% Medium Length Professional CV (Trey Hunner, LaTeXTemplates.com)
% Requires resume.cls in the same directory.
% Tailored for: Summer@EPFL 2027 (computer vision, imaging, machine learning)
%%%%%%%%%%%%%%%%%%%%%%%%%%%%%%%%%%%%%%%%%

\documentclass{resume}

\usepackage[left=0.4in,top=0.4in,right=0.4in,bottom=0.4in]{geometry}
\usepackage[T1]{fontenc}
\usepackage[utf8]{inputenc}
\usepackage{charter}
\usepackage[hidelinks]{hyperref}

\name{Osinachi Chukwujama}
\address{Saint-Étienne, France}
\address{\href{mailto:osinachi.chukwujama@gmail.com}{osinachi.chukwujama@gmail.com} \\ +33 6 51 41 94 14}
\address{\href{https://github.com/vicradon}{github.com/vicradon} \\ \href{https://www.linkedin.com/in/chukwujama-osinachi/}{linkedin.com/in/chukwujama-osinachi}}

\begin{document}

%----------------------------------------------------------------------------------------
\begin{rSection}{Profile}
Photonics master's student with a First Class degree in Mechanical Engineering, working at the meeting
point of optics, imaging and machine learning. I have derived and implemented diffusion and
flow-matching models in PyTorch, deployed vision models on embedded GPUs, worked hands-on with
spatial light modulators, and built optics simulations. I am looking for a summer research project in
computer vision, computational imaging or generative modelling.
\end{rSection}

%----------------------------------------------------------------------------------------
\begin{rSection}{Research Interests}
Computational imaging, generative models for vision (diffusion, flow matching), image restoration and
reconstruction, optimal transport in machine learning.
\end{rSection}

%----------------------------------------------------------------------------------------
\begin{rSection}{Education}

{\bf Erasmus Mundus Joint Master, Intelligent Photonics (iPSRS)} \hfill {September 2026 -- 2028}\\
Université Jean Monnet (France), University of Eastern Finland, Université Paris-Est Créteil\\
\textit{Relevant coursework:} Fourier optics, digital image processing, interferometry, probability and signal analysis,
scientific computing with Python.

{\bf Bachelor of Engineering, Mechanical Engineering} \hfill {November 2017 -- August 2023}\\
Federal University of Technology, Owerri (FUTO), Nigeria. CGPA: 4.56/5.00 (First Class),\\
96th percentile in the department; best graduating student in Industrial and Production Engineering.

\end{rSection}

%----------------------------------------------------------------------------------------
\begin{rSection}{Research and Laboratory Experience}

\begin{rSubsection}{Generative Modelling from First Principles}{2026}{Self-directed research}{}
\item Derived and implemented DDPM (forward/reverse process, ELBO, reparameterisation) and OT-conditional flow matching in PyTorch, covering training and sampling end to end.
\item Studied the links between score matching (denoising and implicit variants, NCSN), flow matching and normalising flows; wrote a forward-diffusion noising tool in Go for synthetic data generation.
\item Proposed a comparison of OT-CFM and DDPM on photonics/laser data (TPU Research Cloud application).
\end{rSubsection}

\begin{rSubsection}{Spatial Light Modulator (SLM) Laboratory Practical}{2026}{iPSRS M1 coursework}{Université Jean Monnet}
\item Generated diffraction gratings and Fresnel lens phase patterns in MATLAB and displayed them on an SLM.
\item Observed and analysed the resulting diffraction patterns, and computed the gratings' spatial frequencies from measurements.
\end{rSubsection}

\begin{rSubsection}{Surveillance Drone with Real-Time Object Detection}{August 2023 -- March 2024}{Undergraduate final project and follow-up, supervised by Prof.\ O.O.\ Obiukwu}{MEE, FUTO}
\item Collected and annotated a three-class dataset, fine-tuned YOLOv5 (70\% mAP) and ran live inference on a drone's camera stream via ONNX.
\item Deployed the model on an NVIDIA Jetson Nano with TensorRT, cutting processing time by 44\%.
\end{rSubsection}

\begin{rSubsection}{Optics Playground}{In progress}{Personal project, \href{https://pre-zeemax.vercel.app/}{pre-zeemax.vercel.app}}{}
\item Interactive 2D ray-tracing simulator with finite-size lenses, being extended into a Fourier optics bench.
\end{rSubsection}

\end{rSection}

%----------------------------------------------------------------------------------------
\begin{rSection}{Skills}

\begin{tabular}{ @{} l @{\hspace{4ex}} p{0.72\textwidth} }
\textbf{Machine learning} & PyTorch: DDPM, flow matching, score matching, CNN/ResNet training, transfer learning; ONNX, TensorRT \\
\textbf{Mathematics} & Probability, variational inference (ELBO), optimal transport for flow matching, Fourier analysis \\
\textbf{Programming} & Python (proficient), MATLAB, Go, TypeScript/JavaScript; Git, Linux, CI/CD, automated testing \\
\textbf{Languages} & English: C1 (IELTS 8.0) \\
\end{tabular}

\end{rSection}

%----------------------------------------------------------------------------------------
\begin{rSection}{Work Experience}

\begin{rSubsection}{Dstrct Group}{March 2025 -- May 2026}{Backend Engineer / Technical Lead}{Lagos, Nigeria}
\item Technical lead on an enterprise platform with a typed backend of nearly 1,000 files; worked directly with clients to turn requirements into engineering tasks.
\end{rSubsection}

\begin{rSubsection}{Hotels.ng}{June 2024 -- January 2026}{Software/DevOps Engineer (part-time)}{Lagos, Nigeria}
\item Built a concurrent Python/Celery application distributing AI workloads over RabbitMQ; wrote CI/CD and monitoring tooling.
\end{rSubsection}

\end{rSection}

%----------------------------------------------------------------------------------------
\begin{rSection}{Publications}
\item \textit{Quantum Computing in Artificial Intelligence: A Review of Quantum Machine Learning Algorithms}, Path of Science.\\ \href{http://dx.doi.org/10.22178/pos.117-25}{doi:10.22178/pos.117-25}
\item \textit{Secure and Resilient Industrial IoT Architectures for Smart Manufacturing: A Comprehensive Review}, JERR (2025).\\ \href{https://doi.org/10.9734/jerr/2025/v27i61548}{doi:10.9734/jerr/2025/v27i61548}
\item \textit{Clinical Applications of 3D-Printed Medical Implants}, JOPIR.\\ \href{https://doi.org/10.69613/22dn3m33}{doi:10.69613/22dn3m33}
\end{rSection}

%----------------------------------------------------------------------------------------
\begin{rSection}{Honours, Teaching and Communication}
\item PTDF Scholarship (2018) for academic excellence.
\item \textbf{Microsoft Certified Trainer} and former \textbf{Microsoft Learn Student Ambassador}: trained over 300 students in software engineering and cloud technologies.
\item \textbf{Director of Research, Mechanical Engineering, FUTO}: ran hands-on training in machine learning, microcontrollers and 3D printing for 50+ students and staff. Co-founder of Hack4FUTO.
\item Technical writer at \href{https://vicradon.hashnode.dev/}{vicradon.hashnode.dev}.
\end{rSection}

\end{document}
